\section{开放性谱系：闭源、开放权重、开源}

课程在历史回顾里不仅列了模型名，也强调了“开放性”本身是一条重要谱系。
对研究者来说，开源/开放权重/闭源不仅是发布策略差异，更决定了你能不能检查训练细节、能不能复现实验、能不能理解失败模式。

\begin{table}[H]
\centering
\begin{tabular}{p{3.0cm}p{3.8cm}p{3.9cm}p{3.4cm}}
\toprule
\textbf{类别} & \textbf{能看到什么} & \textbf{典型例子} & \textbf{研究意义} \\
\midrule
闭源 & 只能通过 API 使用 & GPT-4 / Claude / Gemini & 适合应用，但很难复现 \\
开放权重 & 权重 + 部分训练信息 & Llama / DeepSeek & 可分析模型，仍可能缺失数据细节 \\
开源 & 权重、数据、较完整细节 & OLMo / 部分开源项目 & 最适合教学和系统研究 \\
\bottomrule
\end{tabular}
\caption{开放性不是二元状态，而是一条连续谱}
\end{table}

\begin{figure}[H]
\centering
\begin{tikzpicture}[node distance=0.9cm, every node/.style={font=\small, align=center}]
\node[draw, rounded corners, fill=gray!12, minimum width=2.9cm, minimum height=0.8cm] (a) {closed API};
\node[draw, rounded corners, fill=gray!20, right=0.9cm of a, minimum width=2.9cm, minimum height=0.8cm] (b) {open weights};
\node[draw, rounded corners, fill=gray!28, right=0.9cm of b, minimum width=2.9cm, minimum height=0.8cm] (c) {open source};
\draw[-{Latex[length=3mm]}, thick] (a) -- node[above]{more inspectable} (b);
\draw[-{Latex[length=3mm]}, thick] (b) -- node[above]{more reproducible} (c);
\end{tikzpicture}
\caption{开放性越高，越容易把“经验”变成“可验证的机制”}
\end{figure}

\begin{knowledgebox}{为什么开放性对这门课很重要}
因为课程要训练的是“能独立改造系统的人”，而不是只能消费 API 的人。你越能看到训练和数据细节，就越能把直觉和工程决策对应起来。
\end{knowledgebox}

\begin{importantbox}{历史回顾的真正目的}
不是为了背模型名字，而是为了看清技术路线如何从“可读、可写”变成“可用但不可见”，再回到“尽可能可见、可研究”的方向。
\end{importantbox}

